\documentclass[10pt]{ctexart}
\usepackage[a4paper,margin=1in]{geometry}
\usepackage{amsmath,amssymb,mathtools,bm}
\usepackage{booktabs}
\usepackage{enumitem}
\usepackage{microtype}
\usepackage{hyperref}
\setlist[itemize]{leftmargin=1.5em}
\newcommand{\clip}{\operatorname{clip}}
\title{002：Shared-Backbone Step-wise PPO 技术路线}
\author{}
\date{}
\begin{document}
\maketitle
\vspace{-1.2em}

\section*{目标}
本文给出 Step-wise PPO 的确定实现路线。critic 不作为独立模型维护，而是在 actor language model 上额外挂一个 scalar value head。actor 负责生成 action，value head 在每个环境 step 的 action 边界预测一个 state value。

\section*{Value Head 设计}
设 $h_t\in\mathbb{R}^d$ 是 action 生成前边界位置的 hidden state。主版本使用两层 MLP：
\[
v_\psi(h_t)=W_2\,\mathrm{GELU}(W_1h_t+b_1)+b_2.
\]
linear head 作为低成本 ablation 保留。value head 在 FSDP wrapping 前挂载，使其参与 actor optimizer 和 checkpoint；同步到 vLLM/SGLang rollout engine 前，需要过滤包含 \path{value_head} 的参数。

默认 value position 是
\[
b_t = n_t - L_t - 1,
\]
其中 $n_t$ 是 context+response 总长度，$L_t$ 是 response 长度。value 定义为
\[
V_{\theta,\psi}(s_t)=v_\psi(H_t[b_t]).
\]
这个位置不 attend 到生成 token，因此是严格 state-value，而不是 action-conditioned value。

\section*{Rollout 必需字段}
每个 step sample 需要保存以下字段：
\begin{itemize}
\item \path{uid}：初始任务 group id，用于 episode-level GRPO normalization。
\item \path{traj_uid}：trajectory id，用于恢复 episode 内 step 序列。
\item \path{step_idx}：环境 step index，用于按时间顺序计算 GAE。
\item \path{step_sample_uid}：可选字段，用于 batch expansion/filter 后追踪样本。
\end{itemize}
Tensor 侧需要保存 response mask、token-level reward、immediate step reward、old log probability、old scalar state value、token-broadcast advantage 和 scalar step return。

\section*{Forward 路径}
actor forward 支持 \path{return_step_values=True}。启用后，模型返回语言模型输出和 token-wise scalar values；StepPPO actor 再根据 \path{value_position} 取一个 scalar。old-log-probability 阶段记录：
\[
(\log\pi_{\theta_{\mathrm{old}}}, V^{\mathrm{old}}_t).
\]
actor update 阶段重新计算：
\[
(\log\pi_{\theta}, V_t).
\]

\section*{Step-wise Advantage}
对每个 \path{traj_uid}，按 \path{step_idx} 排序后计算：
\[
\delta_t = r_t + \gamma(1-d_t)V^{\mathrm{old}}_{t+1} - V^{\mathrm{old}}_t,
\]
\[
A_t^{\mathrm{step}} = \delta_t + \gamma\lambda(1-d_t)A_{t+1}^{\mathrm{step}},
\qquad
R_t^{\mathrm{step}} = A_t^{\mathrm{step}} + V^{\mathrm{old}}_t.
\]
这些量全部是 environment-step indexed，不是 token indexed。

\section*{Episode Relative Signal}
对 task group $g$ 内 trajectory $i$，定义 episode return $R_i^{\mathrm{epi}}$。GRPO-style signal 为
\[
A_i^{\mathrm{epi}}
=\frac{R_i^{\mathrm{epi}}-\mu_g}{\sigma_g+\epsilon}.
\]
该 scalar 会广播到 trajectory $i$ 的所有 step。

\section*{Final Advantage 版本}
当前 v1 版本为
\[
A_t^{\mathrm{final,v1}}
=\frac{\widehat A_t^{\mathrm{epi}}+w\widehat A_t^{\mathrm{step}}}{1+w}.
\]
v0 base 版本为
\[
A_t^{\mathrm{final,v0}}
=\frac{\widehat A_t^{\mathrm{epi}}+w A_t^{\mathrm{step}}}{1+w}.
\]
二者区别在于 step GAE advantage 是否先做 batch-level whitening。

\section*{Policy 与 Value Loss}
actor 使用 PPO clipped objective。对 step $t$ 中 token $\ell$：
\[
\rho_{t,\ell}(\theta)
=\exp(\log\pi_\theta(y_{t,\ell})-\log\pi_{\theta_{\mathrm{old}}}(y_{t,\ell})).
\]
policy loss 为
\[
\mathcal{L}_{\mathrm{pg}}
=-\mathbb{E}\left[\min\left(\rho_{t,\ell}A_t^{\mathrm{final}},
\clip(\rho_{t,\ell},1-\epsilon,1+\epsilon)A_t^{\mathrm{final}}\right)\right].
\]
value loss 使用 clipped scalar regression：
\[
\widetilde V_t = V_t^{\mathrm{old}} + \clip(V_t-V_t^{\mathrm{old}},-\epsilon_v,\epsilon_v),
\]
\[
\mathcal{L}_{\mathrm v}
=\frac12\mathbb{E}_t\left[\max\left((V_t-R_t^{\mathrm{step}})^2,
(\widetilde V_t-R_t^{\mathrm{step}})^2\right)\right].
\]
总 loss 为
\[
\mathcal{L}=\mathcal{L}_{\mathrm{pg}}+c_v\mathcal{L}_{\mathrm v}+\beta\mathcal{L}_{\mathrm{KL}}-c_H\mathcal{H}.
\]

\section*{验证计划}
\begin{itemize}
\item 检查每个 step sample 都有 \path{uid}、\path{traj_uid}、\path{step_idx}。
\item 检查 batch shuffle 后 GAE 的 temporal ordering 是否正确。
\item 先跑 2GPU smoke，确认日志里出现 \path{state_values}、\path{step_returns}、\path{actor/value_loss}。
\item 再跑 4GPU WebShop 单 seed，与 GiGPO 和 GRPO 对比。
\end{itemize}

\section*{主要风险}
最大风险是 value target 噪声通过 shared backbone 干扰 actor 表征。第二个风险是额外 value forward 带来的运行开销。需要用 detached-backbone、lower-value-coef 和 microbatch ablation 来定位。

\end{document}
